\documentclass[reqno]{amsart}
\usepackage[T1]{fontenc}
\usepackage{amsmath,amssymb,amsthm}
\usepackage{xr-hyper}
\usepackage[hidelinks]{hyperref}
\externaldocument[ps-]{20261125T010500--poincare-series-and-near-orthogonality__workbook}
\externaldocument[ds-]{20261125T004000--denominators-and-sha__workbook}
\externaldocument[ts-]{20261125T004100--tate-twists-and-sha__workbook}
\externaldocument[es-]{20261125T014500--eisenstein-congruences-and-selmer-classes__workbook}
\externaldocument[pc-]{20261125T014700--pseudocharacters-and-the-length-bound__workbook}

% Chapter number; \ref{chap} gives it.
\newcounter{chap}
% Chapter numbers, generated from 20261125T012000--grh-workbook__guide.org
% by scripts/make-chapters.py.  Do not edit: change the order of the
% index and rerun the script (./preview-site.sh and CI do this).
%
% Usage in a sheet, after \newcounter{chap}:
%   % Chapter numbers, generated from 20261125T012000--grh-workbook__guide.org
% by scripts/make-chapters.py.  Do not edit: change the order of the
% index and rerun the script (./preview-site.sh and CI do this).
%
% Usage in a sheet, after \newcounter{chap}:
%   % Chapter numbers, generated from 20261125T012000--grh-workbook__guide.org
% by scripts/make-chapters.py.  Do not edit: change the order of the
% index and rerun the script (./preview-site.sh and CI do this).
%
% Usage in a sheet, after \newcounter{chap}:
%   \input{chapters}\chapnumber{<sheet ID>}
\makeatletter
\providecommand\chapters@set[2]{\@namedef{chapters@#1}{#2}}
\chapters@set{20261125T001000}{1}% zeros-to-ramp-sums
\chapters@set{20261125T001100}{2}% inert-primes-and-estimate-b
\chapters@set{20261125T005000}{3}% galois-fourier-and-averaging
\chapters@set{20261125T005500}{4}% columns-from-a-sphere
\chapters@set{20261125T000000}{5}% slope-method
\chapters@set{20261126T030100}{6}% a-twist-for-each-prime
\chapters@set{20261125T003000}{7}% the-chain-and-the-square-root
\chapters@set{20261125T010000}{8}% rank-lemma-and-zero-estimate
\chapters@set{20261125T010500}{9}% poincare-series-and-near-orthogonality
\chapters@set{20261124T223000}{10}% elliptic-curves-for-the-proof
\chapters@set{20261125T009500}{11}% fock-model-and-theta-intertwiners
\chapters@set{20261124T223100}{12}% complex-multiplication-and-inert-reduction
\chapters@set{20261125T008000}{13}% one-cm-determinant
\chapters@set{20261124T214500}{14}% supersingular-theta-divisibility
\chapters@set{20261124T214600}{15}% the-isotropic-plane
\chapters@set{20261125T013700}{16}% pairings-on-a-formal-group-of-height-two
\chapters@set{20261125T013500}{17}% hecke-nondegeneracy-on-the-tree
\chapters@set{20261126T030200}{18}% the-bound-is-attained
\chapters@set{20261125T009000}{19}% kronecker-limit-and-chowla-selberg
\chapters@set{20261125T002000}{20}% definite-theta-lifts
\chapters@set{20261125T015500}{21}% siegel-weil-in-the-smallest-case
\chapters@set{20261125T002100}{22}% definite-theta-lifts-and-norms
\chapters@set{20261125T016500}{23}% hecke-operators-and-theta-lifts
\chapters@set{20261125T016700}{24}% divisibility-along-the-chain
\chapters@set{20261125T004000}{25}% denominators-and-sha
\chapters@set{20261125T004100}{26}% tate-twists-and-sha
\chapters@set{20261126T020430}{27}% the-minus-part-by-gauss-sums
\chapters@set{20261125T014100}{28}% congruence-modules
\chapters@set{20261125T014500}{29}% eisenstein-congruences-and-selmer-classes
\chapters@set{20261125T014700}{30}% pseudocharacters-and-the-length-bound
\chapters@set{20261125T011000}{31}% galois-representations-black-box
\chapters@set{20261126T030300}{32}% the-data
\chapters@set{20261126T020210}{33}% bessel-poincare-series-and-their-gram-operators
\chapters@set{20261126T030400}{34}% the-seesaw-identity
\chapters@set{20261126T030500}{35}% divisibility-of-the-composite-lift
\chapters@set{20261126T030600}{36}% theorem-a-short-forms
\chapters@set{20261126T030700}{37}% the-determinant-of-a-chain
\chapters@set{20261126T030800}{38}% the-cost-of-one-lift
\chapters@set{20261126T030900}{39}% denominators-of-projections-to-lifts
\chapters@set{20261126T031000}{40}% twists-with-few-congruences
\chapters@set{20261126T031100}{41}% theorem-c-and-the-main-inequality
\providecommand\chapnumber[1]{%
  \@ifundefined{chapters@#1}%
    {\PackageError{chapters}{Sheet #1 is not listed in the index}%
      {List it in 20261125T012000--grh-workbook__guide.org\MessageBreak
       and run scripts/make-chapters.py.}}%
    {\edef\thechap{\@nameuse{chapters@#1}}}}
\makeatother
\chapnumber{<sheet ID>}
\makeatletter
\providecommand\chapters@set[2]{\@namedef{chapters@#1}{#2}}
\chapters@set{20261125T001000}{1}% zeros-to-ramp-sums
\chapters@set{20261125T001100}{2}% inert-primes-and-estimate-b
\chapters@set{20261125T005000}{3}% galois-fourier-and-averaging
\chapters@set{20261125T005500}{4}% columns-from-a-sphere
\chapters@set{20261125T000000}{5}% slope-method
\chapters@set{20261126T030100}{6}% a-twist-for-each-prime
\chapters@set{20261125T003000}{7}% the-chain-and-the-square-root
\chapters@set{20261125T010000}{8}% rank-lemma-and-zero-estimate
\chapters@set{20261125T010500}{9}% poincare-series-and-near-orthogonality
\chapters@set{20261124T223000}{10}% elliptic-curves-for-the-proof
\chapters@set{20261125T009500}{11}% fock-model-and-theta-intertwiners
\chapters@set{20261124T223100}{12}% complex-multiplication-and-inert-reduction
\chapters@set{20261125T008000}{13}% one-cm-determinant
\chapters@set{20261124T214500}{14}% supersingular-theta-divisibility
\chapters@set{20261124T214600}{15}% the-isotropic-plane
\chapters@set{20261125T013700}{16}% pairings-on-a-formal-group-of-height-two
\chapters@set{20261125T013500}{17}% hecke-nondegeneracy-on-the-tree
\chapters@set{20261126T030200}{18}% the-bound-is-attained
\chapters@set{20261125T009000}{19}% kronecker-limit-and-chowla-selberg
\chapters@set{20261125T002000}{20}% definite-theta-lifts
\chapters@set{20261125T015500}{21}% siegel-weil-in-the-smallest-case
\chapters@set{20261125T002100}{22}% definite-theta-lifts-and-norms
\chapters@set{20261125T016500}{23}% hecke-operators-and-theta-lifts
\chapters@set{20261125T016700}{24}% divisibility-along-the-chain
\chapters@set{20261125T004000}{25}% denominators-and-sha
\chapters@set{20261125T004100}{26}% tate-twists-and-sha
\chapters@set{20261126T020430}{27}% the-minus-part-by-gauss-sums
\chapters@set{20261125T014100}{28}% congruence-modules
\chapters@set{20261125T014500}{29}% eisenstein-congruences-and-selmer-classes
\chapters@set{20261125T014700}{30}% pseudocharacters-and-the-length-bound
\chapters@set{20261125T011000}{31}% galois-representations-black-box
\chapters@set{20261126T030300}{32}% the-data
\chapters@set{20261126T020210}{33}% bessel-poincare-series-and-their-gram-operators
\chapters@set{20261126T030400}{34}% the-seesaw-identity
\chapters@set{20261126T030500}{35}% divisibility-of-the-composite-lift
\chapters@set{20261126T030600}{36}% theorem-a-short-forms
\chapters@set{20261126T030700}{37}% the-determinant-of-a-chain
\chapters@set{20261126T030800}{38}% the-cost-of-one-lift
\chapters@set{20261126T030900}{39}% denominators-of-projections-to-lifts
\chapters@set{20261126T031000}{40}% twists-with-few-congruences
\chapters@set{20261126T031100}{41}% theorem-c-and-the-main-inequality
\providecommand\chapnumber[1]{%
  \@ifundefined{chapters@#1}%
    {\PackageError{chapters}{Sheet #1 is not listed in the index}%
      {List it in 20261125T012000--grh-workbook__guide.org\MessageBreak
       and run scripts/make-chapters.py.}}%
    {\edef\thechap{\@nameuse{chapters@#1}}}}
\makeatother
\chapnumber{<sheet ID>}
\makeatletter
\providecommand\chapters@set[2]{\@namedef{chapters@#1}{#2}}
\chapters@set{20261125T001000}{1}% zeros-to-ramp-sums
\chapters@set{20261125T001100}{2}% inert-primes-and-estimate-b
\chapters@set{20261125T005000}{3}% galois-fourier-and-averaging
\chapters@set{20261125T005500}{4}% columns-from-a-sphere
\chapters@set{20261125T000000}{5}% slope-method
\chapters@set{20261126T030100}{6}% a-twist-for-each-prime
\chapters@set{20261125T003000}{7}% the-chain-and-the-square-root
\chapters@set{20261125T010000}{8}% rank-lemma-and-zero-estimate
\chapters@set{20261125T010500}{9}% poincare-series-and-near-orthogonality
\chapters@set{20261124T223000}{10}% elliptic-curves-for-the-proof
\chapters@set{20261125T009500}{11}% fock-model-and-theta-intertwiners
\chapters@set{20261124T223100}{12}% complex-multiplication-and-inert-reduction
\chapters@set{20261125T008000}{13}% one-cm-determinant
\chapters@set{20261124T214500}{14}% supersingular-theta-divisibility
\chapters@set{20261124T214600}{15}% the-isotropic-plane
\chapters@set{20261125T013700}{16}% pairings-on-a-formal-group-of-height-two
\chapters@set{20261125T013500}{17}% hecke-nondegeneracy-on-the-tree
\chapters@set{20261126T030200}{18}% the-bound-is-attained
\chapters@set{20261125T009000}{19}% kronecker-limit-and-chowla-selberg
\chapters@set{20261125T002000}{20}% definite-theta-lifts
\chapters@set{20261125T015500}{21}% siegel-weil-in-the-smallest-case
\chapters@set{20261125T002100}{22}% definite-theta-lifts-and-norms
\chapters@set{20261125T016500}{23}% hecke-operators-and-theta-lifts
\chapters@set{20261125T016700}{24}% divisibility-along-the-chain
\chapters@set{20261125T004000}{25}% denominators-and-sha
\chapters@set{20261125T004100}{26}% tate-twists-and-sha
\chapters@set{20261126T020430}{27}% the-minus-part-by-gauss-sums
\chapters@set{20261125T014100}{28}% congruence-modules
\chapters@set{20261125T014500}{29}% eisenstein-congruences-and-selmer-classes
\chapters@set{20261125T014700}{30}% pseudocharacters-and-the-length-bound
\chapters@set{20261125T011000}{31}% galois-representations-black-box
\chapters@set{20261126T030300}{32}% the-data
\chapters@set{20261126T020210}{33}% bessel-poincare-series-and-their-gram-operators
\chapters@set{20261126T030400}{34}% the-seesaw-identity
\chapters@set{20261126T030500}{35}% divisibility-of-the-composite-lift
\chapters@set{20261126T030600}{36}% theorem-a-short-forms
\chapters@set{20261126T030700}{37}% the-determinant-of-a-chain
\chapters@set{20261126T030800}{38}% the-cost-of-one-lift
\chapters@set{20261126T030900}{39}% denominators-of-projections-to-lifts
\chapters@set{20261126T031000}{40}% twists-with-few-congruences
\chapters@set{20261126T031100}{41}% theorem-c-and-the-main-inequality
\providecommand\chapnumber[1]{%
  \@ifundefined{chapters@#1}%
    {\PackageError{chapters}{Sheet #1 is not listed in the index}%
      {List it in 20261125T012000--grh-workbook__guide.org\MessageBreak
       and run scripts/make-chapters.py.}}%
    {\edef\thechap{\@nameuse{chapters@#1}}}}
\makeatother
\chapnumber{20261125T014100}
\renewcommand{\thesection}{\thechap.\arabic{section}}

\theoremstyle{definition}
\newtheorem{problem}{Problem}[section]
\theoremstyle{plain}
\newtheorem{theorem}{Theorem}[section]

\newcommand{\Z}{\mathbb{Z}}
\newcommand{\Q}{\mathbb{Q}}
\newcommand{\R}{\mathbb{R}}
\newcommand{\C}{\mathbb{C}}
\newcommand{\F}{\mathbb{F}}
\newcommand{\Oc}{\mathcal{O}}
\newcommand{\T}{\mathbf{T}}
\newcommand{\cc}{\mathfrak{c}}
\newcommand{\pp}{\mathfrak{p}}
\newcommand{\len}{\operatorname{len}}
\newcommand{\Sym}{\operatorname{Sym}}
\newcommand{\Ad}{\operatorname{Ad}}
\newcommand{\Cg}{\mathrm{C}}
\newcommand{\ii}{\mathrm{i}}

\setlength{\emergencystretch}{2em}
\title{Congruence modules}
\author{}
\date{November 25, 2026}

\begin{document}
\maketitle
\refstepcounter{chap}\label{chap}

\paragraph{Prerequisites.}
Modules over a Dedekind domain and the elementary divisor theorem.  Ideals in quadratic fields.  Modular forms of level \(1\): the structure of \(M_k\) and \(S_k\) \cite[Ch.~VII, \S3.2, Thm.~4]{SeA}, Eisenstein series and \(\Delta\) \cite[Ch.~VII, \S\S4.2--4.3]{SeA}, Hecke operators and their action on \(q\)-expansions, and normalized eigenforms \cite[Ch.~VII, \S\S5.3--5.5]{SeA}; Serre writes the weight as \(2k\) and numbers the Bernoulli numbers differently.  The Petersson product, for which the Hecke operators are self-adjoint \cite[\S5.5]{DS}.  From Chapter~\ref{ds-chap}, \emph{Denominators and the Stickelberger element}: the cost of a denominator (Problem~\ref{ds-pr:enlarge}) and Theorem~\ref{ds-thm:HR}.

In Chapter~\ref{ds-chap} a linear map that does not respect integral structures cost \(\log [M + T(L) : M]\).  One step of the proof asks for an operator clearing the denominators of an orthogonal projection onto the image of a theta lift restricted to good forms.  It tests that operator on a restricted set of integral forms.  Projection onto a sum of whole Hecke eigenspaces is a model for this problem.  What does such a projection cost, and what controls the cost?  On the full lattice the answer is a congruence module.  We compute it for level \(1\), in weight \(12\), where it is \(\Z/691\Z\), and in weight \(24\), where it detects a congruence modulo the unique prime ideal above \(144169\) between the two newforms.

\section{Congruence modules}\label{sec:modules}

Let \(\Oc\) be a Dedekind domain with fraction field \(K\).  A \emph{lattice} in a finite-dimensional \(K\)-vector space \(V\) is a finitely generated \(\Oc\)-submodule that spans \(V\).  Let \(V = V_1 \oplus V_2\), let \(\pi_1, \pi_2\) be the two projections, and let \(\Lambda \subset V\) be a lattice.  The \emph{congruence module} of \(V_1\) in \(\Lambda\) is
\[
  \Cg(\Lambda; V_1) := \pi_1(\Lambda) / (\Lambda \cap V_1) .
\]

\begin{problem}\label{pr:symmetry}
\(\pi_1(\Lambda) = V_1 \cap (\Lambda + V_2)\), and
\[
  \Cg(\Lambda; V_1) \cong \Lambda / \bigl((\Lambda \cap V_1) \oplus (\Lambda \cap V_2)\bigr) \cong \Cg(\Lambda; V_2) .
\]
In particular \(\Cg(\Lambda; V_1)\) is a finitely generated torsion \(\Oc\)-module, and it vanishes exactly when \(\Lambda = (\Lambda \cap V_1) \oplus (\Lambda \cap V_2)\).
\end{problem}

\begin{problem}\label{pr:denominator}
\((\Lambda + \pi_1(\Lambda))/\Lambda \cong \Cg(\Lambda; V_1)\).  For \(\Oc = \Z\), the least positive integer \(c\) with \(c\,\pi_1(\Lambda) \subseteq \Lambda\) is the exponent of \(\Cg(\Lambda; V_1)\).
\end{problem}

So for \(\Oc = \Z\) the index \([\Lambda + \pi_1(\Lambda) : \Lambda]\) is \(\lvert \Cg(\Lambda; V_1) \rvert\): measured like the index cost of Problem~\ref{ds-pr:enlarge} (stated there for injective maps), the projection costs \(\log \lvert \Cg(\Lambda; V_1) \rvert\).

\begin{problem}\label{pr:toy}
Let \(N \ge 1\), \(V = \Q^2 = \Q e_1 \oplus \Q e_2\), and \(\Lambda = \{(a,b) \in \Z^2 : a \equiv b \pmod N\}\).  Then \(\Cg(\Lambda; \Q e_1) \cong \Z/N\Z\), whereas \(\Cg(\Z^2; \Q e_1) = 0\).
\end{problem}

The lattices of interest carry an action of a Hecke algebra.  Let \(X\) be a finite set, and let \(\T \subseteq \prod_{x \in X} \Oc\) be an \(\Oc\)-subalgebra, finitely generated as an \(\Oc\)-module, which separates the points of \(X\).  Write \(t(x)\) for the coordinates of \(t \in \T\).  Let \(\T\) act on \(V = \bigoplus_{x \in X} V_x\) by \(t \cdot v = (t(x) v_x)_x\), and let \(\Lambda\) be a \(\T\)-stable lattice.  The \emph{congruence ideal} of \(x\) is
\[
  \cc_x := \{ t(x) : t \in \T,\ t(y) = 0 \text{ for all } y \ne x \} \subseteq \Oc ,
\]
and we write \(\Cg(\Lambda; V_x)\) for the congruence module with respect to \(V = V_x \oplus \bigoplus_{y \ne x} V_y\).

\begin{problem}\label{pr:annihilate}
\(\cc_x\) annihilates \(\Cg(\Lambda; V_x)\).  If \(\Oc\) is a discrete valuation ring with normalized valuation \(v\), then \(\len_\Oc \Cg(\Lambda; V_x) \le \dim V_x \cdot \min v(\cc_x)\).
\end{problem}

\begin{problem}\label{pr:free}
If \(\Lambda \cong \T\) as a \(\T\)-module, then \(\Cg(\Lambda; V_x) \cong \Oc/\cc_x\).  For \(\T = \{(a,b) \in \Z^2 : a \equiv b \pmod N\}\) acting on itself, \(\cc_x = N\Z\) at both points.
\end{problem}

Thus \(\cc_x\) annihilates the obstruction to integral projection onto \(V_x\).  When the lattice is free of rank one over \(\T\), this obstruction module is isomorphic to \(\Oc/\cc_x\).

A congruence module can also be read off from a pairing.  Suppose that \(\Lambda\) carries a \emph{perfect} pairing \(\langle\,,\rangle \colon \Lambda \times \Lambda \to \Oc\), that is, an \(\Oc\)-bilinear form for which \(\lambda \mapsto \langle \lambda, \cdot \rangle\) is an isomorphism \(\Lambda \to \operatorname{Hom}_\Oc(\Lambda, \Oc)\), and that \(\langle V_1, V_2 \rangle = \langle V_2, V_1 \rangle = 0\).  For a lattice \(L \subset V_1\) let \(L^\vee := \{ v \in V_1 : \langle v, L \rangle \subseteq \Oc \}\).

\begin{problem}\label{pr:selfdual}
\(\pi_1(\Lambda) = (\Lambda \cap V_1)^\vee\).  Hence \(\Cg(\Lambda; V_1) \cong (\Lambda \cap V_1)^\vee / (\Lambda \cap V_1)\).
\end{problem}

\begin{problem}\label{pr:gram}
In the setting of Problem~\ref{pr:selfdual}:
(i) if \(f\) is a basis of \(\Lambda\cap V_1\) over \(\Oc\), then \(\Cg(\Lambda; V_1) \cong \Oc/\langle f, f \rangle\Oc\);
(ii) if \(x,y\) form an \(\Oc\)-basis of \(\Lambda\cap V_1\) with \(\langle x, x \rangle = \langle y, y \rangle = 0\), then \(\Cg(\Lambda; V_1) \cong (\Oc/\langle x, y \rangle\Oc) \oplus (\Oc/\langle y, x \rangle\Oc)\).
For \(\Lambda = \Z^2\) with the dot product and \(V_1 = \Q(1,1)\), \(V_2 = \Q(1,-1)\), the congruence module is \(\Z/2\Z\).
\end{problem}

For a self-dual lattice the congruence module of an orthogonal summand is its discriminant module, the quotient in Problem~\ref{pr:selfdual}.  In Section~\ref{sec:hida} the discriminant of the cohomological eigenlattice detects the primes in the support of the eigenform's congruence module in the range of Hida's theorem.  Compare \cite[\S4.4, Lemma~4.17]{DDT}, where a power of the congruence ideal of the Hecke algebra is shown to lie in the discriminant ideal of the eigenform's piece of a self-dual lattice, with equality when the lattice is free over the Hecke algebra.

\section{The Eisenstein line in weight \texorpdfstring{\(12\)}{12}}\label{sec:691}

For even \(k \ge 4\) let \(M_k\) and \(S_k\) be the spaces of modular forms and cusp forms of weight \(k\) for \(\mathrm{SL}_2(\Z)\), identified with their \(q\)-expansions, and let \(M_k(\Z) := M_k \cap \Z[[q]]\) and \(S_k(\Z) := S_k \cap \Z[[q]]\).  Let \(\sigma_{k-1}(n) := \sum_{d \mid n} d^{k-1}\) and
\[
  E_k := 1 - \frac{2k}{B_k} \sum_{n \ge 1} \sigma_{k-1}(n) q^n , \qquad \Delta := q \prod_{n \ge 1} (1 - q^n)^{24} = \sum_{n \ge 1} \tau(n) q^n ,
\]
so that \(E_4 = 1 + 240 \sum \sigma_3(n) q^n\), \(E_6 = 1 - 504 \sum \sigma_5(n) q^n\), and \(B_{12} = -691/2730\) \cite[Ch.~VII, \S4]{SeA}.  Write \(M_k(\Q):=M_k\cap\Q[[q]]\) and \(S_k(\Q):=S_k\cap\Q[[q]]\).  We use the decomposition \(M_k(\Q)=\Q E_k\oplus S_k(\Q)\), in which both summands are stable under the Hecke operators.

\begin{problem}\label{pr:basis12}
\(M_{12}(\Z) = \Z E_4^3 \oplus \Z \Delta\).
\end{problem}

\begin{problem}\label{pr:691}
\(E_4^3 = E_{12} + \frac{432000}{691}\Delta\).  Hence \(\pi_\Delta(M_{12}(\Z)) = \frac{1}{691}\Z\Delta\), and \(\Cg(M_{12}(\Z); \Q\Delta) \cong \Z/691\Z\).
\end{problem}

The same computation for the Eisenstein line works in every weight.

\begin{problem}\label{pr:bernoulli}
Let \(k \ge 4\) be even, and let \(N_k\) be the numerator of \(B_k/2k\) in lowest terms.  Then \(\Cg(M_k(\Z); \Q E_k) \cong \Z/N_k\Z\).  For \(k = 4, 6, 8, 10, 14\) one has \(N_k = \pm1\), and \(\lvert N_{16} \rvert = 3617\), \(\lvert N_{22} \rvert = 77683 = 131 \cdot 593\), \(\lvert N_{24} \rvert = 236364091 = 103 \cdot 2294797\).
\end{problem}

Let \(\T_{12}\) be the \(\Z\)-algebra of endomorphisms of \(M_{12}(\Q)\) generated by the Hecke operators \(T_n\), \(n \ge 1\), normalized so that \(T_n \Delta = \tau(n)\Delta\) and \(T_n E_{12} = \sigma_{11}(n) E_{12}\) \cite[Ch.~VII, \S5]{SeA}.  Map \(\T_{12}\) to \(\Z \times \Z\) by its eigenvalues on \((\Delta, E_{12})\).

\begin{problem}\label{pr:hecke12}
Each \(T_n\) preserves \(M_{12}(\Z)\), and \(\T_{12} = \{(a,b) \in \Z^2 : a \equiv b \pmod{691}\}\).  Thus the congruence ideal of \(\Delta\) is \(691\Z\).
\end{problem}

After localization at \(691\), the length bound of Problem~\ref{pr:annihilate} is attained.  Consider now the global \(\T_{12}\)-module.

\begin{problem}\label{pr:notfree}
\(M_{12}(\Z)\) is not free of rank one over \(\T_{12}\).
\end{problem}  The prime \(691\) is also the one in Theorem~\ref{ds-thm:HR}: since \(691 \mid B_{12}\), the eigenspace \(A^{(679)}\) of the \(691\)-part of the class group of \(\Q(\zeta_{691})\) is nonzero.  That the same prime measures both is not an accident; Chapter~\ref{es-chap} builds the class from the congruence.

\section{Two newforms of weight \texorpdfstring{\(24\)}{24}}\label{sec:24}

The space \(S_{24}\) has dimension \(2\) \cite[Ch.~VII, \S3.2]{SeA}, and its two normalized eigenforms are not rational.  Let \(D := 144169\), a prime with \(D \equiv 1 \pmod{24}\), let \(K := \Q(\sqrt D)\), \(\Oc := \Oc_K\), and \(\pp := \sqrt D\,\Oc\).

\begin{problem}\label{pr:basis24}
\(S_{24}(\Z) = \Z g_1 \oplus \Z g_2\), where \(g_1 = \Delta E_4^3 - 696\Delta^2 = q + 195660 q^3 + \cdots\) and \(g_2 = \Delta^2 = q^2 - 48q^3 + \cdots\).
\end{problem}

On \(S_{24}\) one has \(a_n(T_2 f) = a_{2n}(f) + 2^{23} a_{n/2}(f)\), with \(a_{n/2} := 0\) for odd \(n\).

\begin{problem}\label{pr:T2}
\(T_2 g_1 = 20468736\, g_2\) and \(T_2 g_2 = g_1 + 1080\, g_2\).  The characteristic polynomial of \(T_2\) on \(S_{24}\) is \(x^2 - 1080x - 20468736\), with discriminant \(24^2 D\) and roots \(\lambda_\pm = 540 \pm 12\sqrt D\).  The normalized eigenforms are \(f_\pm = g_1 + \lambda_\pm g_2\).
\end{problem}

Let \(\T_{24}\) be the \(\Z\)-algebra generated by the \(T_n\) on \(S_{24}\).

\begin{problem}\label{pr:T24}
\(\T_{24} = \Z[T_2] = \Z \oplus \Z T_2\), and \(S_{24}(\Z) = \T_{24}\, g_2\) is free of rank one over \(\T_{24}\).
\end{problem}

Let \(\Lambda := S_{24}(\Z) \otimes \Oc\) in \(V := S_{24}(\Z) \otimes K = K f_+ \oplus K f_-\), and let \(\T_{24} \otimes \Oc\) act on it, through its eigenvalues, as a subalgebra of \(\Oc \times \Oc\).

\begin{problem}\label{pr:cong24}
The congruence ideal of \(f_+\) is \((\lambda_+ - \lambda_-)\Oc = 24\sqrt D\,\Oc\).  Hence \(\Cg(\Lambda; K f_+) \cong \Oc/24\sqrt D\,\Oc\).  Verify the module isomorphism also by direct projection.
\end{problem}

\begin{problem}\label{pr:modp}
\(a_n(f_+) \equiv a_n(f_-) \pmod{24\sqrt D\,\Oc}\) for all \(n\).  The primes \(2\) and \(3\) split in \(K\), and \(\Cg(\Lambda; K f_+)\) has length \(1\) at \(\pp\), length \(3\) at each prime above \(2\), and length \(1\) at each prime above \(3\).
\end{problem}

The congruence modulo \(\pp\), the ramified prime of the Hecke field, is the interesting one.  The primes \(2\) and \(3\) are smaller than the weight, outside the range \(p > k\) of Hida's theorem below (Theorem~\ref{thm:hida}).

\section{Hida's identity}\label{sec:hida}

Which primes occur in the congruence module of a cuspidal eigenform?  Let \(k \ge 12\) be even, let \(f = \sum a(n) q^n \in S_k\) be a normalized eigenform, let \(K \subset \C\) be a number field containing the coefficients of all normalized eigenforms in \(S_k\), \(\Oc := \Oc_K\), and \(\Lambda := S_k(\Z) \otimes \Oc\).  Let \(Y \subset S_k(\Z) \otimes K\) be the span of the other normalized eigenforms, and write \(\Cg(f) := \Cg(\Lambda; K f)\) for the decomposition \(S_k(\Z) \otimes K = K f \oplus Y\).  Let
\[
  \langle g, h \rangle := \int_{\mathrm{SL}_2(\Z)\backslash \mathcal H} g(z)\,\overline{h(z)}\, y^{k} \,\frac{dx\,dy}{y^2}
\]
be the Petersson product.

\begin{problem}\label{pr:orth}
\(Y\) is the orthogonal complement of \(f\), and the projection to \(K f\) is \(g \mapsto \frac{\langle g, f \rangle}{\langle f, f \rangle} f\).  Hence \(\Cg(f) \cong J_f/\Oc\), where \(J_f\) is the \(\Oc\)-module generated by the numbers \(\langle g, f \rangle / \langle f, f \rangle\) with \(g \in S_k(\Z)\).  For \(f = f_+\) of Section~\ref{sec:24}, \(J_f = (24\sqrt D)^{-1}\Oc\).
\end{problem}

The projection divides by \(\langle f,f\rangle\).  Hida's theorem detects the primes above \(p>k\) at which the projection fails to preserve the lattice.  For \(p\) prime write \(a(p) = \alpha_p + \beta_p\) with \(\alpha_p\beta_p = p^{k-1}\), and let
\[
  L(\Sym^2 f, s) := \prod_p \bigl(1 - \alpha_p^2 p^{-s}\bigr)^{-1} \bigl(1 - p^{k-1-s}\bigr)^{-1} \bigl(1 - \beta_p^2 p^{-s}\bigr)^{-1} \qquad (\operatorname{Re} s > k) .
\]
In the unitary normalization, the \emph{adjoint \(L\)-function} is \(L(s, \Ad f) := L(\Sym^2 f, s + k - 1)\); its value at \(s = 1\) is \(L(\Sym^2 f, k)\).

\begin{problem}\label{pr:euler}
Let \(\alpha,\beta\in\C\), \(w:=\alpha\beta\), and \(a_m:=\sum_{i=0}^m\alpha^i\beta^{m-i}\).  Then, as a formal power series in \(X\),
\[
  \sum_{m \ge 0} a_m^2 X^m = \frac{1 + wX}{(1 - \alpha^2 X)(1 - wX)(1 - \beta^2 X)} .
\]
Hence, for \(\operatorname{Re}s>k\), \(\sum_{n \ge 1} a(n)^2 n^{-s} = \zeta(s-k+1)\, L(\Sym^2 f, s) / \zeta(2s-2k+2)\).
\end{problem}

Rankin's method \cite{Ra} computes the residue of this Dirichlet series at \(s = k\); in terms of \(L(\Sym^2 f, s)\), which is holomorphic there \cite{Sh},
\begin{equation}\label{eq:rankin}
  \langle f, f \rangle = \frac{(k-1)!}{2^{2k-1} \pi^{k+1}} \, L(\Sym^2 f, k)
\end{equation}
\cite[p.~109, (3)--(5)]{Za}.  For \(f = \Delta\) both sides equal \(1.0353620568\ldots \times 10^{-6}\), the value of \(\lVert\Delta\rVert^2\) found from Petersson's formula in Section~\ref{ps-sec:petersson} of Chapter~\ref{ps-chap}.

The transcendental part of \(\langle f, f \rangle\) is measured by periods.  Let
\[
  r_f(X) := \int_0^{\ii\infty} f(z) (X - z)^{k-2} \, dz ,
\]
and let \(r_f^+\) and \(r_f^-\) be its even and odd parts in \(X\).  By Manin's theorem there are \(\omega_f^\pm \in \C^\times\) with \(r_f^\pm \in \omega_f^\pm K[X]\) \cite{Mn}.  For \(f = \Delta\) one has \cite{KZ}
\begin{gather*}
  r_\Delta^+ = \omega^+ \Bigl( \tfrac{36}{691}(X^{10} - 1) + X^2 - 3X^4 + 3X^6 - X^8 \Bigr), \\
  r_\Delta^- = \omega^- \bigl( 4X^9 - 25X^7 + 42X^5 - 25X^3 + 4X \bigr) .
\end{gather*}

\begin{problem}\label{pr:relation}
The two polynomials in parentheses, and \(X^{10} - 1\), satisfy \(P(X) + X^{10} P(-1/X) = 0\).
\end{problem}

The polynomial \(X^{10} - 1\) is a coboundary in the Eichler--Shimura cohomology, and the prime in front of it is the \(691\) of Problem~\ref{pr:691}.  For a cusp form one discards this direction.  Fix a prime \(\pp\) of \(\Oc\) above a prime \(p > k\), let \(\Oc_\pp\) be the localization of \(\Oc\) at \(\pp\), and let \(\Omega_f^\pm \in \C^\times\) be Hida's periods, defined by generators of the \(f\)-eigenlines in the parabolic cohomology of \(\mathrm{SL}_2(\Z)\) with coefficients in polynomials of degree at most \(k - 2\) over \(\Oc_\pp\); the period polynomial represents the class of \(f\) there, modulo the coboundary \(X^{k-2} - 1\).  For the numerical checks below we use instead the displayed period-polynomial coefficients: \(\Omega_f^\pm\) such that the coefficients of \(X^n\) in \(r_f^\pm\) with \(0 < n < k-2\) lie in \(\Omega_f^\pm \Oc_\pp\) and generate it.  Identifying this normalization with Hida's in general requires an integral comparison of period polynomials and cohomology classes modulo \(X^{k-2} - 1\).  The following theorem is due to Hida \cite{Hi1,Hi2}, with a contribution of Ribet \cite{Ri}; see \cite[Thm.~1]{Gh}.  We use it as a black box.

\begin{theorem}[Hida]\label{thm:hida}
The number \(\eta_f := \langle f, f \rangle / (\ii\,\Omega_f^+ \Omega_f^-)\) lies in \(K\).  If \(\pp\) lies above a prime \(p > k\), then \(\Cg(f) \otimes \Oc_\pp \ne 0\) if and only if \(\eta_f \in \pp\Oc_\pp\).
\end{theorem}

By \eqref{eq:rankin}, \(\eta_f\) is, up to a \(\pp\)-adic unit, the algebraic part \(L(1, \Ad f)/(\pi^{k+1}\ii\,\Omega_f^+\Omega_f^-)\) of the adjoint \(L\)-value.  In outline, the proof applies Problem~\ref{pr:gram}(ii) to the parabolic cohomology lattice, self-dual for \(p > k\), whose \(f\)-part is spanned by isotropic \(\pm\) lines; the main difficulty is to compare this cohomological congruence module with \(\Cg(f)\) \cite[\S3]{Gh}; in weight \(2\) see \cite[\S4.4]{DDT}.  Two checks with PARI/GP \cite{PARI}.  For \(f=\Delta\), take the coefficient of \(X^2\) in \(r_\Delta^+\) as \(\Omega^+\) and take \(\Omega^-=\omega^-\).  The computation gives \(\eta_\Delta=\pm 2^{-10}\), in agreement with the absence of cuspidal congruence primes for \(S_{12}=\C\Delta\).  Replacing \(\Omega^+\) by \((36/691)\Omega^+\) would give \(v_{691}(\eta_\Delta)=1\).  For \(f=f_+\) of Section~\ref{sec:24}, take \(K=\Q(\sqrt D)\) and use the coefficients of \(X^2\) in \(r_f^+\) and of \(X\) in \(r_f^-\) for \(\Omega^+\) and \(\Omega^-\), respectively.  These coefficients generate the respective fractional ideals of the interior coefficients at every prime above \(p>24\).  The norm \(N_{K/\Q}(\eta_f)\) is \(D\) times a rational number whose numerator and denominator have no prime factor greater than \(24\).  Thus \(v_\pp(\eta_f)=1\) at \(\pp=\sqrt D\,\Oc\), as in Problem~\ref{pr:modp}.

\paragraph{Comparison with the full lattice.}
For a Hecke-stable lattice \(\Lambda\) as in Section~\ref{sec:modules}, let \(V_1=\bigoplus_{x\in I}V_x\) for a nonempty subset \(I\) of \(X\), and project along the sum of the remaining eigenspaces.  On the full lattice the obstruction is \(\Cg(\Lambda;V_1)\), by Problems~\ref{pr:symmetry} and~\ref{pr:denominator}, and the ideal \(\bigcap_{x\in I}\cc_x\) annihilates it, as in Problem~\ref{pr:annihilate}.  For a single cuspidal eigenline, Theorem~\ref{thm:hida} detects its congruence primes above \(p>k\) by an adjoint \(L\)-value.  For the Eisenstein congruence of \(\Delta\) at \(691\), Ribet's construction gives an unramified quotient of order \(691\) of \(A^{(679)}\) (Chapter~\ref{es-chap}).  For odd primes \(\ell\), Chapter~\ref{ts-chap} gives the Selmer length over \(\Q\) for positive odd twists \(e\) and identifies the classes killed by \(\ell\) with homomorphisms on a class-group eigenspace when \((\ell-1)\nmid e\).  Chapter~\ref{es-chap} constructs a class from an Eisenstein congruence, and Problem~\ref{pc-pr:sha-lower} of Chapter~\ref{pc-chap} bounds the congruence length of one eigenform by a Selmer length.

\paragraph{The step this chapter serves.}
The chapter serves one step of the proof: clearing the denominators of an orthogonal projection onto a space of theta lifts.  The corresponding statement of the Lean proof \cite{GL} is \texttt{ClearsProjection}, in \texttt{grh/src/New/Determinant/Twist/FitTwists.lean}; the corresponding bound is \texttt{denominatorLe\_clearingBound}, in \texttt{grh/src/New/Congruences/FirstOrder/DenominatorBound.lean}, with its specialization \texttt{denominatorLe\_range\_goodLift}, in \texttt{grh/src/New/Congruences/FitTwistsExist.lean}.

\paragraph{Exercise.}
For \(k = 16, 18, 20, 22, 26\) the space \(S_k\) is spanned by one normalized eigenform \(f_k\), and \(a_\ell(f_k) \equiv 1 + \ell^{k-1} \pmod{N_k}\) for every prime \(\ell\).

\section*{Hints}

\subsection*{\S\ref{sec:modules}}
\begin{description}
\item[\ref{pr:symmetry}] Consider \(\Lambda \to \pi_1(\Lambda)/(\Lambda \cap V_1)\), \(\lambda \mapsto \pi_1(\lambda)\), and compute its kernel.
\item[\ref{pr:denominator}] The map \(\pi_1(\Lambda) \to (\Lambda + \pi_1(\Lambda))/\Lambda\).
\item[\ref{pr:toy}] \(\pi_1(\Lambda) = \Z e_1\).
\item[\ref{pr:annihilate}] Apply to \(\Lambda\) an element of \(\T\) that vanishes off \(x\).  For the length, \(\pi_x(\Lambda)\) is free of rank \(\dim V_x\).
\item[\ref{pr:free}] Identify \(V_x\) with \(K\) by \(t \mapsto t(x)\); then \(\pi_x(\T) = \Oc\).
\item[\ref{pr:selfdual}] \(\Lambda \cap V_1\) is a direct summand of \(\Lambda\), so an \(\Oc\)-valued functional on it extends to \(\Lambda\); then use that the pairing is perfect.
\item[\ref{pr:gram}] Compute the dual basis.
\end{description}

\subsection*{\S\ref{sec:691}}
\begin{description}
\item[\ref{pr:basis12}] Subtract a multiple of \(E_4^3\) to kill the constant term; \(S_{12} = \C\Delta\).
\item[\ref{pr:691}] Compare the coefficients of \(1\) and \(q\).  Then use \(\gcd(432000, 691) = 1\).
\item[\ref{pr:bernoulli}] The projection to \(\Q E_k\) is \(g \mapsto a_0(g)E_k\), and some \(E_4^aE_6^b\) lies in \(M_k(\Z)\) with constant term \(1\).
\item[\ref{pr:hecke12}] For the first claim use the formula for \(a_n(T_m f)\) in \cite[Ch.~VII, \S5.3]{SeA}.  For the second, Problem~\ref{pr:691} gives one inclusion; for the other use \(T_2\), \(T_3\) and \(\gcd(2073, 176896) = 691\).  Here \(\tau(2)=-24\) and \(\tau(3)=252\).
\item[\ref{pr:notfree}] For \(g = aE_4^3 + b\Delta\), compute the index of \(\T_{12}\,g\) in \(M_{12}(\Z)\) using \((1,1)\) and \((691,0)\); note \(432000 \equiv 125 \pmod{691}\).
\end{description}

\subsection*{\S\ref{sec:24}}
\begin{description}
\item[\ref{pr:basis24}] As in Problem~\ref{pr:basis12}, with \(S_{24} = \Delta M_{12}\).
\item[\ref{pr:T2}] By Problem~\ref{pr:basis24} the first two coefficients of \(A g_1 + B g_2\) are \(A, B\); a normalized eigenform has \(a_1 = 1\) and \(a_2\) = its \(T_2\)-eigenvalue \cite[Ch.~VII, \S5.4]{SeA}.  One has \(a_4(g_1) = 12080128\) and \(a_4(g_2) = 1080\).
\item[\ref{pr:T24}] \(\T_{24} \otimes \Q = \Q[T_2]\), and the matrix of \(T_2\) in the basis \(g_1, g_2\) has an entry \(1\).
\item[\ref{pr:cong24}] Write \(t = a + bT_2\) with \(a, b \in \Oc\), and apply Problem~\ref{pr:free}.  For the direct check, write \(g_1, g_2\) in terms of \(f_\pm\).
\item[\ref{pr:modp}] Write \(f_+ - f_-\) in the basis \(g_1, g_2\).  For the splitting, reduce \(D\) modulo \(8\) and \(3\).
\end{description}

\subsection*{\S\ref{sec:hida}}
\begin{description}
\item[\ref{pr:orth}] Use simultaneous Hecke eigenspaces.  For the integral lattice use a basis with \(a_i(g_j)=\delta_{ij}\) in the first \(\dim S_k\) coefficients, and recall \(a_1(f)=1\).  For the example use Problem~\ref{pr:cong24}.
\item[\ref{pr:euler}] First take \(\alpha\ne\beta\) and use the difference of powers: expand \(a_m^2\) as a combination of \(\alpha^{2m}\), \(w^m\) and \(\beta^{2m}\), and sum three geometric series.  Both sides of the resulting coefficient identities are polynomials in \(\alpha,\beta\).
\end{description}

\begin{thebibliography}{DDT}
\bibitem[DDT]{DDT} H.~Darmon, F.~Diamond and R.~Taylor, \emph{Fermat's Last Theorem}, Current Developments in Mathematics \textbf{1995} (1995), no.~1, 1--154.
\bibitem[DS]{DS} F.~Diamond and J.~Shurman, \emph{A First Course in Modular Forms}, Grad.\ Texts in Math.\ 228, Springer, 2005.
\bibitem[Gh]{Gh} E.~Ghate, \emph{An introduction to congruences between modular forms}, in: Current Trends in Number Theory (Allahabad, 2000), Hindustan Book Agency, New Delhi, 2002, 39--58.
\bibitem[GL]{GL} L.~Sernau, \emph{GRH for Dirichlet \(L\)-functions, in Lean~4}, a shorter rewriting of the Lean proof, text by Claude, \url{https://github.com/lsernau/grh-lean}, commit \texttt{72fead1}, readable edition \texttt{grh/}.
\bibitem[Hi1]{Hi1} H.~Hida, \emph{Congruences of cusp forms and special values of their zeta functions}, Invent.\ Math.\ \textbf{63} (1981), 225--261.
\bibitem[Hi2]{Hi2} H.~Hida, \emph{On congruence divisors of cusp forms as factors of the special values of their zeta functions}, Invent.\ Math.\ \textbf{64} (1981), 221--262.
\bibitem[KZ]{KZ} W.~Kohnen and D.~Zagier, \emph{Modular forms with rational periods}, in: Modular Forms (Durham, 1983), R.~A.~Rankin (ed.), Ellis Horwood, Chichester, 1984, 197--249.
\bibitem[Mn]{Mn} Yu.~I.~Manin, \emph{Periods of parabolic forms and \(p\)-adic Hecke series}, Math.\ USSR-Sb.\ \textbf{21} (1973), 371--393.
\bibitem[PARI]{PARI} The PARI~Group, \emph{PARI/GP}, version 2.17.1, Univ.\ Bordeaux, \url{https://pari.math.u-bordeaux.fr/}.
\bibitem[Ra]{Ra} R.~A.~Rankin, \emph{Contributions to the theory of Ramanujan's function \(\tau(n)\) and similar arithmetical functions.~I, II}, Proc.\ Cambridge Philos.\ Soc.\ \textbf{35} (1939), 351--356 and 357--372.
\bibitem[Ri]{Ri} K.~A.~Ribet, \emph{Mod \(p\) Hecke operators and congruences between modular forms}, Invent.\ Math.\ \textbf{71} (1983), 193--205.
\bibitem[SeA]{SeA} J.-P.~Serre, \emph{A Course in Arithmetic}, Grad.\ Texts in Math.\ 7, Springer, 1973.
\bibitem[Sh]{Sh} G.~Shimura, \emph{On the holomorphy of certain Dirichlet series}, Proc.\ London Math.\ Soc.\ (3) \textbf{31} (1975), 79--98.
\bibitem[Za]{Za} D.~Zagier, \emph{Modular forms whose Fourier coefficients involve zeta-functions of quadratic fields}, in: Modular Functions of One Variable VI, Lecture Notes in Math.\ 627, Springer, 1977, 105--169.
\end{thebibliography}

\end{document}
